Pour chacune des affirmations suivantes, indiquer si elle est vraie ou fausse en justifiant votre choix.
Une réponse non argumentée ne sera pas prise en compte dans l'évaluation.

\begin{enumerate}
\item On considère la fonction $f$ définie sur l'intervalle $]1~;+\infty[$ par :
\[
f(x)=\dfrac{x-1}{\sqrt{x^2-1}}
\]

Affirmation 1 : La fonction $f$ admet pour limite $1$ en $+\infty$.

\item On considère la suite $\left(w_n\right)$ définie par $w_0 = 1$ et pour tout entier naturel $n$ par :
\[
w_{n+1}=w_n+2n+3
\]

Affirmation 2 : Pour tout entier naturel $n$,
\[
w_n=(n+1)^2
\]

\item Soit $p$ un nombre réel tel que $0 < p < 1$.

On considère une variable aléatoire $X$ qui suit la loi binomiale de paramètres $3$ et $p$.

On note $P\left(X=1\right)$ la probabilité de l'évènement $(X=1)$.

Affirmation 3 : $P\left(X=1\right)=3p-6p^2+3p^3$.

\item On considère la suite $\left(v_n\right)$ définie pour tout entier naturel $n > 1$ par :
\[
v_n=\int_0^1 \e^{nx}\,\d x
\]

Affirmation 4 : Pour tout entier naturel $n > 1$,
\[
v_n=\dfrac{\e^n}{n}
\]

\item On colorie en rouge, jaune ou noir chacune des $16$ cases d'un quadrillage.

Affirmation 5 : On peut réaliser $\binom{16}{3}$ coloriages différents.
\end{enumerate}
